\chapter*{Note on Translation and Transliteration}
\addcontentsline{toc}{chapter}{Note on Translation and Transliteration}
\markboth{Translation and Transliteration}{Translation and Transliteration}

This book draws on three distinct linguistic traditions --- classical
Arabic, English-language scientific and philosophical literature, and
the technical terminology of modern mathematics. Coordination among
them requires explicit conventions, which are stated here.

\section*{Arabic transliteration}

The book uses a simplified transliteration system rather than the
full \textit{International Journal of Middle East Studies} (IJMES)
scheme. The simplification trades scholarly precision for readability:
specialists who require IJMES-level precision can recover it from the
glossary, where each term is given in both transliterated and full
diacritical form.

The simplifications are as follows. Long vowels are not marked with a
macron in the running text (so we write \arb{rizq} rather than
\arb{r\bar{\i}zq}, where appropriate). The letter \arb{ʿayn} is
transliterated with an apostrophe (\arb{'}) in the running text and
with the IJMES symbol \arb{ʿ} in the glossary. The \arb{hamza}
is similarly simplified to an apostrophe. Emphatic consonants are not
distinguished from their non-emphatic counterparts in the running
text; the glossary entries provide the full transliteration for
specialists.

The exceptions are proper names of persons widely recognized in
English. \textit{Ibn Arabi} appears as such, not as
\textit{Ibn al-`Arab\bar{\i}}. \textit{Mulla Sadra} appears in this
common form rather than \textit{Mull\bar{a} \d{S}adr\bar{a}}.

Arabic terms are set in italic on first use and on later occurrences
where the term is being thematized. After the first few occurrences,
common terms (\arb{rizq}, \arb{barakah}, \arb{tawhid},
\arb{niyyah}, \arb{qadar}, \arb{wujud}, \arb{tajalli},
\arb{waqf}, \arb{zakat}, \arb{riba}) are typically set in
roman type.

\section*{Quranic citations}

Quranic verses are cited in the form ``Quran X:Y'' where X is the
\textit{surah} number and Y is the verse number within the surah. So
``Quran 17:31'' refers to verse 31 of \textit{Surat al-Isra'} (Surah
17). The number Y always refers to the verse in the standard
Egyptian (\textit{Hafs `an `Asim}) numbering, which differs slightly
from the older Flugel numbering used in some nineteenth-century
European editions.

English translations of Quranic verses are taken, where not otherwise
specified, from \textit{The Clear Quran} by Mustafa Khattab (2015).
Where alternative translations are needed for technical reasons, they
are noted at the point of use. Where the author has departed from
Khattab's rendering, the departure is indicated and explained.

\section*{Hadith citations}

Hadith are cited in two ways. Major collections are cited by the
collection name and the standard volume/book/hadith numbering. So
``\hadith{Bukhari} 1385'' refers to hadith number 1385 in the standard
edition of the \textit{Sahih al-Bukhari}. ``\hadith{Tirmidhi} 1653''
refers to hadith 1653 in the standard edition of the
\textit{Jami` al-Tirmidhi}. The standard editions are those of the
1995 Riyadh print runs unless otherwise noted.

When a hadith's grading (sahih, hasan, da'if, etc.) is relevant to the
argument, the grading is given. Following standard scholarly
practice, the grading is the assessment of the major hadith critics
(al-Albani, Ibn Hajar, etc.) where one is widely accepted; where
gradings are disputed, both are noted.

\section*{Classical Islamic sources}

When citing classical works, the title is given in transliteration
(italicized), then in English translation (in parentheses, on first
use), then by section or chapter where applicable. So:
``Ibn Arabi, \textit{al-Futuhat al-Makkiyyah} (\textit{The Meccan
Revelations}), Chapter 198.'' Subsequent citations use a short form:
``\textit{Futuhat}, Ch. 73.''

Page references are provided where the edition matters; many of these
texts have been printed in multiple editions with different
paginations. Where possible, the chapter and section numbering ---
which is consistent across editions --- is used.

\section*{Mathematical and scientific notation}

Standard conventions are used throughout. Where a notation is
non-standard or where multiple conventions exist, the convention used
in this book is stated explicitly in the Notation section that
follows. The book follows the conventions of the relevant field:
Hilbert space notation follows the Dirac convention used in physics
texts (\(|\psi\rangle, \langle\phi|\)); probability notation follows
the conventions of measure-theoretic probability (Billingsley,
\textit{Probability and Measure}, 3rd ed.); economic notation follows
Mas-Colell, Whinston, and Green (\textit{Microeconomic Theory}).

Where notations from different fields conflict (the symbol \(P\)
denotes both probability and price; the symbol \(S\) denotes both
entropy and supply), the meaning is stated explicitly at first use in
each chapter. Context generally disambiguates after that.

\section*{Translations of technical terms}

A persistent challenge in this book is that several Islamic technical
terms have no precise English equivalent. The pragmatic conventions
adopted are these:

\begin{itemize}[leftmargin=2em]
    \item For terms with no acceptable English equivalent
    (\arb{barakah}, \arb{tajalli}, \arb{wujud},
    \arb{tawhid}, \arb{insan al-kamil}), the Arabic is used
    throughout, italicized on first use, with a definition given at
    that point and again in the glossary.
    \item For terms with a near-equivalent that distorts meaning
    (\arb{rizq} as ``sustenance,'' \arb{niyyah} as
    ``intention,'' \arb{qadar} as ``decree''), both the Arabic
    and the near-equivalent are used. The Arabic is preferred where
    technical precision is required; the English where the gloss is
    sufficient.
    \item For terms with acceptable English equivalents (\arb{`ilm}
    as ``knowledge,'' \arb{`adl} as ``justice''), the English is
    used in the running text and the Arabic supplied in the glossary.
\end{itemize}

A list of all Arabic terms used in the book, with their definitions
and transliteration conventions, appears in the Glossary in the back
matter and is summarized for quick reference in the section
``Permitted Use of Arabic Terms'' that follows the Notation section.

\cleardoublepage
